\section{Задание}

В рамках выполнения лабораторной работы \textit{по варианту №629486} необходимо выполнить следующий перечень заданий:
\begin{enumerate}
  \item Для своей программы из \href{https://github.com/aaalvaaas/web3}{лабораторной работы №3} по дисциплине «Веб-программирование» реализовать:
  \begin{itemize}
    \item MBean, считающий общее число установленных пользователем точек, а также число точек, не попадающих в область. В случае, если количество установленных пользователем точек стало кратно 15, разработанный MBean должен отправлять оповещение об этом событии.
    \item MBean, определяющий процентное отношение «промахов» к общему числу кликов пользователя по координатной плоскости.
  \end{itemize}
  \item С помощью утилиты JConsole провести мониторинг программы:
  \begin{itemize}
    \item Снять показания MBean-классов, разработанных в ходе выполнения задания 1.
    \item Определить наименование и версию JVM, поставщика виртуальной машины Java и номер её сборки.
  \end{itemize}
  \item С помощью утилиты VisualVM провести мониторинг и профилирование программы:
  \begin{itemize}
    \item Снять график изменения показаний MBean-классов, разработанных в ходе выполнения задания 1, с течением времени.
    \item Определить имя класса, объекты которого занимают наибольший объём памяти JVM; определить пользовательский класс, в экземплярах которого находятся эти объекты.
  \end{itemize}
  \item С помощью утилиты VisualVM и профилировщика IDE NetBeans, Eclipse или Idea локализовать и устранить проблемы с производительностью в \href{https://github.com/itmo-fse-ld-2026/4-performance/tree/main/http-unit}{программе}. По результатам локализации и устранения проблемы необходимо составить отчёт, в котором должна содержаться следующая информация:
  \begin{itemize}
    \item Описание выявленной проблемы.
    \item Описание путей устранения выявленной проблемы.
    \item Подробное \textit{(со скриншотами)} описание алгоритма действий, который позволил выявить и локализовать проблему.
  \end{itemize}
\end{enumerate}

Студент должен обеспечить возможность воспроизведения процесса поиска и локализации проблемы по требованию преподавателя.